\label{sec:interior-exterior-boundary}

\begin{topicbox}{Inside, Outside, and Edge}
\begin{exposition}
Relative to a set $E$, each point of the line is interior to $E$, exterior to $E$, or on its boundary. These three regions partition $\mathbb{R}$ and refine the open/closed distinction.
\end{exposition}
\end{topicbox}
\begin{definitionbox}{Definition (Interior Point)}
\begin{definition}[Interior Point]
\label{def:interior-point}
Let $A\subseteq\mathbb{R}$ and $x\in\mathbb{R}$.  The point $x$ is an
interior point of $A$ if there exists $r>0$ such that
\[
  B_r(x)\subseteq A.
\]
Equivalently,
\[
  x\text{ is an interior point of }A
  \Longleftrightarrow
  (\exists r>0)(\forall y\in\mathbb{R})
  \bigl(|x-y|<r\Longrightarrow y\in A\bigr).
\]
\end{definition}
\end{definitionbox}
\begin{remark*}[Standard quantified statement]
\[
x \text{ is an interior point of } E \Longleftrightarrow \exists r>0\;(B_{r}(x)\subseteq E).
\]
\end{remark*}
\begin{remark*}[Interpretation]
An interior point of $E$ sits inside $E$ with room to spare: a whole ball around it lies in $E$.
\end{remark*}
\begin{dependencies}
\begin{itemize}
  \item \hyperref[def:open-ball]{Open Ball on the Real Line}
\end{itemize}
\end{dependencies}
\begin{definitionbox}{Definition (Exterior Point)}
\begin{definition}[Exterior Point]
\label{def:exterior-point}
Let $A\subseteq\mathbb{R}$ and $x\in\mathbb{R}$.  The point $x$ is an
exterior point of $A$ if there exists $r>0$ such that
\[
  B_r(x)\subseteq\mathbb{R}\setminus A.
\]
Equivalently,
\[
  x\text{ is an exterior point of }A
  \Longleftrightarrow
  (\exists r>0)(B_r(x)\cap A=\varnothing).
\]
\end{definition}
\end{definitionbox}
\begin{remark*}[Standard quantified statement]
\[
x \text{ is an exterior point of } E \Longleftrightarrow \exists r>0\;(B_{r}(x)\subseteq \mathbb{R}\setminus E).
\]
\end{remark*}
\begin{remark*}[Interpretation]
An exterior point of $E$ is interior to the complement: a whole ball around it misses $E$.
\end{remark*}
\begin{dependencies}
\begin{itemize}
  \item \hyperref[def:open-ball]{Open Ball on the Real Line}
\end{itemize}
\end{dependencies}
\begin{definitionbox}{Definition (Boundary Point)}
\begin{definition}[Boundary Point]
\label{def:boundary-point}
Let $A\subseteq\mathbb{R}$ and $x\in\mathbb{R}$.  The point $x$ is a
boundary point of $A$ if every open ball centered at $x$ meets both $A$ and
its complement; that is,
\[
  (\forall r>0)(B_r(x)\cap A\neq\varnothing)
\]
and
\[
  (\forall r>0)(B_r(x)\cap(\mathbb{R}\setminus A)\neq\varnothing).
\]
\end{definition}
\end{definitionbox}
\begin{remark*}[Standard quantified statement]
\[
x \text{ is a boundary point of } E \Longleftrightarrow \forall r>0\;(B_{r}(x)\cap E\neq\emptyset \;\land\; B_{r}(x)\cap(\mathbb{R}\setminus E)\neq\emptyset).
\]
\end{remark*}
\begin{remark*}[Interpretation]
A boundary point of $E$ is approached from both sides: every ball around it meets both $E$ and its complement.
\end{remark*}
\begin{dependencies}
\begin{itemize}
  \item \hyperref[def:open-ball]{Open Ball on the Real Line}
\end{itemize}
\end{dependencies}
\begin{definitionbox}{Definition (Interior of a Set)}
\begin{definition}[Interior of a Set]
\label{def:interior-of-set}
Let $A\subseteq\mathbb{R}$.  The \emph{interior} of $A$ is the set of all
interior points of $A$:
\[
  \operatorname{int}(A)
  :=
  \{x\in\mathbb{R}:(\exists r>0)(B_r(x)\subseteq A)\}.
\]
Equivalently,
\[
  x\in\operatorname{int}(A)
  \Longleftrightarrow
  (\exists r>0)(\forall y\in\mathbb{R})
  \bigl(|x-y|<r\Longrightarrow y\in A\bigr).
\]
\end{definition}
\end{definitionbox}
\begin{remark*}[Standard quantified statement]
\[
\operatorname{int}(E)=\{\,x\in\mathbb{R}: \exists r>0\;(B_r(x)\subseteq E)\,\}.
\]
\end{remark*}
\begin{remark*}[Interpretation]
The interior is the largest open set inside $E$; $E$ is open exactly when $E=\operatorname{int}(E)$. The boundary is what $E$ and its complement share, and a set is closed exactly when it contains its boundary.
\end{remark*}
\begin{dependencies}
\begin{itemize}
  \item \hyperref[def:interior-point]{Interior Point}
\end{itemize}
\end{dependencies}
